% thqg_introduction_supplement_body.tex
% Extended introduction: from bar complex to holographic datum,
% shadow connection, gravitational complexity, duality and phase space,
% collision to Yang--Baxter, perturbative finiteness, and guide.

% ======================================================================
%  SECTION 2: FROM BAR COMPLEX TO HOLOGRAPHIC DATUM
% ======================================================================

\section{From bar complex to holographic datum}
\label{sec:thqg-intro-bar-to-holographic}
\index{holographic datum!from bar complex}
\index{Swiss-cheese operad!in holographic datum}

The holographic triangle of \S\ref{sec:holographic-triangle} was
presented as a diagram of three vertices and three edges.  The following
constructs it from the bar complex alone, by progressively extracting
structure from the identity $D_\cA^2 = 0$.  The extraction proceeds
in four stages: the Swiss-cheese factorization that separates the
bar complex into a bulk direction and a boundary direction; the
bulk-boundary-line triangle that identifies each vertex categorically;
the holographic datum that packages the triangle into a single
Maurer--Cartan element; and the reconstruction theorems that recover
the triangle from the datum.

\subsection{The Swiss-cheese factorization}
\label{subsec:thqg-intro-swiss-cheese}
\index{Swiss-cheese operad!factorization}
\index{bar complex!as Swiss-cheese coalgebra}

Let $\cA$ be a chiral algebra on a smooth projective curve~$X$.  The
bar complex $\barB_X(\cA)$ is a dg coalgebra on $\operatorname{Ran}(X)$,
and its differential~$d_{\barB}$ acts by residue extraction on
$\overline{C}_n(X)$.  The first structural observation is that this
differential decomposes into two commuting pieces when the curve~$X$
is embedded in a surface.

Let $\Sigma = X \times \mathbb{R}$, with $X$ the holomorphic direction
and $\mathbb{R}$ the topological direction.  The Fulton--MacPherson
compactification $\overline{C}_n(\Sigma)$ fibers over the ordered
configuration space $\operatorname{Conf}_n(\mathbb{R})$, and the bar
differential decomposes as
\begin{equation}\label{eq:thqg-intro-swiss-cheese-decomposition}
d_{\barB,\Sigma}
\;=\;
d_C + d_R,
\end{equation}
where $d_C$ acts in the $X$-direction (extracting OPE residues along
holomorphic collision divisors) and $d_R$ acts in the $\mathbb{R}$-direction
(extracting residues along the topological ordering).  The identity
$d_{\barB,\Sigma}^2 = 0$ decomposes as
\begin{equation}\label{eq:thqg-intro-swiss-cheese-squares}
d_C^2 = 0, \qquad d_R^2 = 0, \qquad \{d_C, d_R\} = 0.
\end{equation}
The first equation is the bar nilpotence of~$\cA$ on~$X$.  The second
is the coassociativity of the bar coproduct.  The third is the
compatibility of the OPE with the ordered product, the defining
condition of a Swiss-cheese coalgebra structure.

\begin{definition}[Swiss-cheese structure on the bar complex]
\label{def:thqg-intro-swiss-cheese-structure}
\index{Swiss-cheese coalgebra!on bar complex}
The \emph{Swiss-cheese structure} on $\barB_X(\cA)$ is the
pair $(d_C, d_R)$ satisfying~\eqref{eq:thqg-intro-swiss-cheese-squares}.
It makes $\barB_X(\cA)$ an algebra over the Swiss-cheese operad
$\mathrm{SC}^{\mathrm{ch,top}}$ of
Theorem~\ref{thm:bar-swiss-cheese}: the closed sector
$d_C$~is $\Einf$-chiral (factorization on~$X$), the open sector
$d_R$~is $\Eone$-topological (ordered configurations
on~$\mathbb{R}$), and the mixed structure is the bar differential
restricted to configurations where some points collide holomorphically
while others collide topologically.
\end{definition}

The Swiss-cheese decomposition is the geometric origin of the
bulk-boundary-line triangle.  The closed sector produces the bulk:
factorization homology $\int_X \cA$ in the $X$-direction.  The open
sector produces the lines: modules over the Koszul dual.  The
boundary is the interface between the two: the chiral algebra~$\cA$
itself, which lives at the $\mathbb{R} = 0$ slice of~$\Sigma$.

\subsection{The bulk-boundary-line triangle}
\label{subsec:thqg-intro-bbl-triangle}
\index{bulk-boundary-line triangle!construction}

The three vertices of the holographic triangle
(\S\ref{subsec:triangle-diagram}) are now constructed as follows.

\smallskip\noindent\textbf{Boundary.}
The boundary algebra is $\cA$ itself.  It carries the full
OPE data and is the input to the bar construction.  The closed-sector
differential $d_C$ acts on $\barB_X(\cA)$ and produces the bar
cohomology $H^*(\barB_X(\cA))$, from which the Koszul dual is
extracted.

\smallskip\noindent\textbf{Bulk.}
The bulk is the ambient complex
$\mathbf{C}_g(\cA) = R\Gamma(\overline{\mathcal{M}}_g,\,
\mathcal{Z}(\cA))$, where
$\mathcal{Z}(\cA) = R\mathcal{H}om_{\cA\text{-bimod}}(\cA, \cA)$
is the center local system.  The ambient complex receives a map
from the bar complex via the trace:
\begin{equation}\label{eq:thqg-intro-bar-to-ambient}
\operatorname{tr}_{g} \colon
\barB^{(g)}_X(\cA)
\;\longrightarrow\;
\mathbf{C}_g(\cA),
\end{equation}
which extracts the central part of each bar element by contracting
all internal indices.  The trace map intertwines $d_{\barB}$ with the
differential on the ambient complex; the image
$\mathbf{Q}_g(\cA) = \operatorname{im}(\operatorname{tr}_g)$ is one
half of the complementarity decomposition
(Theorem~\ref{thm:quantum-complementarity-main}).

\smallskip\noindent\textbf{Lines.}
The line-operator category is $\mathcal{C}(\cA) = \cA^!\text{-}\mathsf{mod}$
(\S\ref{subsec:vertex-lines}).  Its objects are
$\cA^!$-modules; its morphisms are chain maps.  The open-sector
differential $d_R$ on $\barB_X(\cA)$ is the module structure: for
$M \in \cA^!\text{-}\mathsf{mod}$, the action of $\cA^!$ on~$M$ is
encoded by the ordered-configuration structure of the bar complex
restricted to the topological direction.

The three vertices form a commutative triangle
\eqref{eq:supp-holographic-triangle-diagram} whose commutativity is
Theorem~A (the Verdier intertwining).  The triangle is not built from
three independent objects connected by ad hoc functors; it is a single
object (the bar complex with its Swiss-cheese structure) viewed
from three angles.

\subsection{Edges of the triangle}
\label{subsec:thqg-intro-triangle-edges}
\index{holographic triangle!edges}

The three edges connecting the vertices are:

\begin{enumerate}[label=\textup{(\roman*)}]
\item \emph{Boundary $\to$ Bulk} (the bar map).
  The bar construction $\barB_X \colon \cA \mapsto \barB_X(\cA)$
  maps the boundary algebra to the bar coalgebra.  The trace
  map~\eqref{eq:thqg-intro-bar-to-ambient} projects the bar
  coalgebra into the ambient complex.  At genus~$0$, the bar map
  is an embedding; at genus $g \ge 1$, its image acquires curvature
  $\kappa(\cA) \cdot \omega_g$.

\item \emph{Bulk $\to$ Lines} (the holographic cofiber).
  The complementarity decomposition
  $\mathbf{C}_g(\cA) \simeq \mathbf{Q}_g(\cA) \oplus \mathbf{Q}_g(\cA^!)$
  identifies the cofiber of the boundary embedding:
  $\operatorname{cofib}(\mathbf{Q}_g(\cA) \hookrightarrow \mathbf{C}_g(\cA))
  \simeq \mathbf{Q}_g(\cA^!)$.
  The cofiber is the dual shadow, which controls the line-operator
  category via the center-module adjunction.

\item \emph{Boundary $\to$ Lines} (Koszul duality).
  The bottom edge of the triangle is the composite
  $\cA \mapsto \barB_X(\cA) \mapsto H^*(\barB_X(\cA)) \mapsto
  (H^*(\barB_X(\cA)))^\vee = \cA^!$.
  On the Koszul locus, this is bar-dualize: the bar construction
  followed by Verdier duality.  The module category
  $\cA^!\text{-}\mathsf{mod}$ is the categorical output of this
  composite.
\end{enumerate}

\subsection{The holographic datum from the MC element}
\label{subsec:thqg-intro-holographic-from-mc}
\index{holographic modular Koszul datum!construction from MC}

The holographic datum
$\mathcal{H}(\cA) = (\cA, \cA^!, \mathcal{C}, r(z), \Theta_\cA,
\nabla^{\mathrm{hol}})$
of~\eqref{eq:supp-holographic-datum-intro} is not defined by listing
six independent objects.  It is extracted from the single
Maurer--Cartan element $\Theta_\cA$ by a sequence of projections.

\begin{enumerate}[label=\textup{(\roman*)}]
\item The boundary algebra~$\cA$ is the input.
\item The Koszul dual~$\cA^!$ is the graded linear dual of the bar
  cohomology coalgebra: $\cA^! = (H^*(\barB_X(\cA)))^\vee$.  By Theorem~A,
  $\mathbb{D}_{\operatorname{Ran}}\,\barB_X(\cA) \simeq \cA^!_\infty$
  (the homotopy Koszul dual algebra, whose underlying complex is
  equivalent to $\barB_X(\cA^!)$).
\item The line-operator category is $\mathcal{C} = \cA^!\text{-}\mathsf{mod}$.
\item The spectral $r$-matrix is the binary genus-$0$ collision residue
  of~$\Theta_\cA$:
  \begin{equation}\label{eq:thqg-intro-r-matrix-extraction}
  r(z) \;=\; \operatorname{Res}^{\mathrm{coll}}_{0,2}(\Theta_\cA)
  \;\in\;
  \cA^! \,\widehat{\otimes}\, \cA^![\![z^{-1}, z]\!].
  \end{equation}
  This extraction is the content of the collision filtration
  (\S\ref{subsec:thqg-intro-collision-definition}).
\item The MC element $\Theta_\cA$ itself: the bar-intrinsic
  construction $\Theta_\cA := D_\cA - d_\cA^{(0)}$
  (Theorem~\ref{thm:mc2-bar-intrinsic}).
\item The shadow connection $\nabla^{\mathrm{hol}}$ is constructed
  from $\Theta_\cA$ by the modular shadow projection
  (\S\ref{sec:thqg-intro-shadow-connection}).
\end{enumerate}

The datum is subordinated to the single identity $D_\cA^2 = 0$.  Every
entry is a projection of~$\Theta_\cA$; every consistency condition
between the entries is a consequence of the MC equation
$D\Theta + \frac{1}{2}[\Theta,\Theta] = 0$.  The holographic
programme of this monograph is the systematic extraction of physics
from this identity.

\subsection{Reconstruction from the datum}
\label{subsec:thqg-intro-reconstruction}
\index{holographic datum!reconstruction theorems}

The holographic datum determines the triangle, and the triangle
determines the datum.  The reconstruction theorems are:

\begin{enumerate}[label=\textup{(\roman*)}]
\item \emph{Boundary from bulk.}
  On the Koszul locus,
  $\cA \simeq \Omega_X(\mathbb{D}_{\operatorname{Ran}}\,
  \mathbf{Q}_g(\cA^!))$:
  the boundary algebra is recovered from the dual shadow by the cobar
  construction applied to the Verdier dual of the bulk
  (Theorem~\ref{thm:higher-genus-inversion}).

\item \emph{Bulk from boundary.}
  The ambient complex is reconstructed by factorization homology:
  $\mathbf{C}_g(\cA) \simeq \int_{\overline{\mathcal{M}}_g}
  \mathcal{Z}(\cA)$, where the center~$\mathcal{Z}(\cA)$ is a
  computable bimodule functor of~$\cA$.

\item \emph{Lines from $r$-matrix.}
  The meromorphic tensor structure on $\cA^!\text{-}\mathsf{mod}$
  is reconstructed from the $r$-matrix $r(z)$: the fusion product
  $M \otimes_{r(z)} N$ is determined by the spectral data of~$r(z)$,
  and the Yang--Baxter equation for~$r(z)$ guarantees associativity
  of the tensor product.

\item \emph{MC element from the datum.}
  Conversely, the six-fold datum determines $\Theta_\cA$ uniquely:
  the boundary~$\cA$ determines the bar complex, the bar complex
  determines~$D_\cA$, and $\Theta_\cA = D_\cA - d_\cA^{(0)}$ is
  the positive-genus correction.  The entries $\cA^!$, $\mathcal{C}$,
  $r(z)$, and $\nabla^{\mathrm{hol}}$ are then projections.
\end{enumerate}

The full content of the reconstruction is that the six-fold datum is
a \emph{complete invariant} of the modular homotopy type
(Proposition~\ref{prop:modular-homotopy-classification}): two chiral
algebras have equivalent holographic data if and only if their modular
bar constructions are quasi-isomorphic.

\subsection{The four recovery theorems}
\label{subsec:thqg-intro-four-recovery}
\index{recovery theorems!holographic}

The following four theorems make the holographic datum rigid.  Each
recovers one piece of the datum from another, using only the MC
equation.

\begin{theorem}[Collision residue $=$ twisting morphism;
\ClaimStatusProvedElsewhere]
\label{thm:thqg-intro-collision-twisting}
The binary genus-$0$ collision residue of~$\Theta_\cA$ is the
universal twisting morphism
$\tau \colon \barB_X(\cA) \to \cA$ restricted to bar degree~$2$:
\begin{equation}\label{eq:thqg-intro-collision-equals-twisting}
\operatorname{Res}^{\mathrm{coll}}_{0,2}(\Theta_\cA)
\;=\;
\tau|_{\barB^2_X(\cA)}.
\end{equation}
The proof is Theorem~\textup{\ref{thm:collision-residue-twisting}}.
\end{theorem}

\begin{theorem}[Arnold $\Rightarrow$ CYBE;
\ClaimStatusProvedElsewhere]
\label{thm:thqg-intro-arnold-cybe}
The classical Yang--Baxter equation for~$r(z)$ is a formal consequence
of the Arnold relation on~$\overline{C}_3(X)$.  The three-term sum
\[
[r_{12}(z_1 - z_2),\, r_{13}(z_1 - z_3)]
+ [r_{12}(z_1 - z_2),\, r_{23}(z_2 - z_3)]
+ [r_{13}(z_1 - z_3),\, r_{23}(z_2 - z_3)]
\;=\; 0
\]
follows from
$\eta_{12} \wedge \eta_{23} + \eta_{23} \wedge \eta_{31}
+ \eta_{31} \wedge \eta_{12} = 0$
applied to the arity-$3$ MC equation.
The proof is Theorem~\textup{\ref{thm:collision-depth-2-ybe}}.
\end{theorem}

\begin{theorem}[Shadow/KZ comparison on the affine Kac--Moody surface;
\ClaimStatusProvedElsewhere]
\label{thm:thqg-intro-shadow-kz}
For $\cA = \widehat{\mathfrak g}_k$ on the affine Kac--Moody
comparison surface, the shadow connection $\nabla^{\mathrm{hol}}$
restricted to genus~$0$ with~$n$ marked points identifies with the
Knizhnik--Zamolodchikov connection on the bundle of conformal blocks:
\[
\nabla^{\mathrm{KZ}}
\;=\;
d - \sum_{i < j} \frac{r_{ij}}{z_i - z_j}\,dz_i
\;=\;
\nabla^{\mathrm{hol}}_{0,n}\big|_{\text{genus $0$}}.
\]
The proof is Theorem~\textup{\ref{thm:shadow-connection-kz}}.
\end{theorem}

\begin{theorem}[Quartic obstruction $=$ $L_\infty$ bracket;
\ClaimStatusProvedElsewhere]
\label{thm:thqg-intro-quartic-linfty}
The quartic resonance class
$\mathfrak{Q}(\cA) \in H^*(\Defcyc^{(4)}(\cA))$ is the
$L_\infty$-obstruction governing the quartic extension problem for the
shadow package, equivalently the obstruction to extending
$\Theta_\cA^{\le 3}$ to $\Theta_\cA^{\le 4}$:
\[
\mathfrak{Q}(\cA)
\;=\;
[\ell_4^{\mathrm{min}}]
\;\in\;
H^2(\Defcyc^{(4)}(\cA)|_{g=0}).
\]
This is the $r = 4$ case of the shadow-Massey identification
\textup{(}Proposition~\textup{\ref{prop:shadow-massey-identification}}\textup{)}.
\end{theorem}

Together, these four theorems show that the holographic datum is not
an external imposition on the bar complex but an intrinsic structure
of the identity $D_\cA^2 = 0$, read at successively finer resolution.


% ======================================================================
%  SECTION 3: THE SHADOW CONNECTION AND CONFORMAL BLOCKS
% ======================================================================

\section{The shadow connection and conformal-block comparison surfaces}
\label{sec:thqg-intro-shadow-connection}
\index{shadow connection|textbf}
\index{conformal blocks!and shadow connection}
\index{Knizhnik--Zamolodchikov connection|textbf}
\index{BPZ equations|textbf}

The modular MC element~$\Theta_\cA$ determines a flat connection on
the moduli of marked curves.  On the general theorematic surface,
this connection governs the transport of the bar-cohomological /
derived-coinvariant package across
$\overline{\mathcal{M}}_{g,n}$.  On benchmark comparison
surfaces, its genus-$0$ affine specialization identifies with the
Knizhnik--Zamolodchikov connection, while on the genus-$0$ Virasoro
degenerate-module surface it gives the Virasoro conformal Ward
connection and its higher collision residues yield the BPZ differential
equations.  Its flatness is guaranteed by the MC equation.

\subsection{Construction of the shadow connection}
\label{subsec:thqg-intro-shadow-connection-construction}
\index{shadow connection!construction}

Let $\pi \colon \overline{\mathcal{C}}_{g,n} \to \overline{\mathcal{M}}_{g,n}$
be the universal curve with~$n$ marked sections
$\sigma_1, \ldots, \sigma_n \colon \overline{\mathcal{M}}_{g,n}
\to \overline{\mathcal{C}}_{g,n}$.  A derived coinvariant /
protected-state package is the sheaf on~$\overline{\mathcal{M}}_{g,n}$
defined by
\begin{equation}\label{eq:thqg-intro-conformal-block-bundle}
\mathcal{V}_{g,n}(\cA; M_1, \ldots, M_n)
\;:=\;
R\pi_*\!\left(
  \cA \;\otimes\;
  \bigotimes_{i=1}^n \sigma_i^*\!(M_i)
\right),
\end{equation}
where $M_1, \ldots, M_n$ are $\cA$-modules inserted at the marked
points.  On comparison surfaces where classical conformal blocks are
defined, the fiber over a smooth curve~$\Sigma_g$ with marked points
$z_1, \ldots, z_n$ is the space of coinvariants:
\[
V_{g,n}(\cA; M_1, \ldots, M_n)
\;=\;
\frac{M_1 \otimes \cdots \otimes M_n}{
  \cA(\Sigma_g \setminus \{z_1, \ldots, z_n\})
  \cdot (M_1 \otimes \cdots \otimes M_n)
},
\]
the quotient by the diagonal action of sections of~$\cA$ regular
away from the marked points.

\begin{definition}[The shadow connection]
\label{def:thqg-intro-shadow-connection}
\index{shadow connection!definition}
The \emph{modular shadow connection} on this derived coinvariant /
protected-state package is the connection
\begin{equation}\label{eq:thqg-intro-shadow-connection-formula}
\nabla^{\mathrm{hol}}_{g,n}
\;=\;
d \;-\; \operatorname{Sh}_{g,n}(\Theta_\cA),
\end{equation}
where $d$ is the de~Rham differential on~$\overline{\mathcal{M}}_{g,n}$
and $\operatorname{Sh}_{g,n}(\Theta_\cA)$ is the genus-$(g,n)$
shadow: the projection of the MC element~$\Theta_\cA$ to the
component acting on the corresponding fiber at genus~$g$
with~$n$ marked points.  On benchmark comparison surfaces, this is
the corresponding action on conformal blocks.

Explicitly, the shadow $\operatorname{Sh}_{g,n}$ is a graph sum over
stable graphs~$\Gamma$ of type $(g,n)$:
\begin{equation}\label{eq:thqg-intro-shadow-graph-sum}
\operatorname{Sh}_{g,n}(\Theta_\cA)
\;=\;
\sum_{\Gamma \in G(g,n)}
\frac{1}{|\operatorname{Aut}(\Gamma)|}
\;\Theta_{\cA,\Gamma},
\end{equation}
where $G(g,n)$ is the set of stable graphs of type $(g,n)$ and
$\Theta_{\cA,\Gamma}$ is the amplitude of~$\Theta_\cA$ on the
graph~$\Gamma$, obtained by placing the multilinear operations
of~$\Theta_\cA$ at the vertices of~$\Gamma$ and contracting along
the internal edges using the bar pairing.
\end{definition}

\subsection{Flatness from the MC equation}
\label{subsec:thqg-intro-flatness}
\index{shadow connection!flatness}
\index{Maurer--Cartan equation!and flatness}

The fundamental property of the shadow connection is flatness:

\begin{proposition}[Flatness of the shadow connection;
\ClaimStatusProvedElsewhere]
\label{prop:thqg-intro-flatness}
The connection $\nabla^{\mathrm{hol}}_{g,n}$ is flat:
\begin{equation}\label{eq:thqg-intro-flatness}
(\nabla^{\mathrm{hol}}_{g,n})^2
\;=\;
0.
\end{equation}
\end{proposition}

\begin{proof}[Proof sketch]
Expand:
\[
(\nabla^{\mathrm{hol}})^2
= d^2 - d\,\operatorname{Sh}(\Theta) - \operatorname{Sh}(\Theta) \wedge d
+ \operatorname{Sh}(\Theta) \wedge \operatorname{Sh}(\Theta)
= -d\,\operatorname{Sh}(\Theta) + \tfrac{1}{2}[\operatorname{Sh}(\Theta), \operatorname{Sh}(\Theta)].
\]
The shadow projection is compatible with the MC equation:
$d\,\operatorname{Sh}(\Theta) = \operatorname{Sh}(D\Theta)$ and
$\operatorname{Sh}(\Theta) \wedge \operatorname{Sh}(\Theta)
= \operatorname{Sh}([\Theta, \Theta])$.  Therefore
\[
(\nabla^{\mathrm{hol}})^2
= -\operatorname{Sh}\bigl(D\Theta + \tfrac{1}{2}[\Theta, \Theta]\bigr)
= 0,
\]
where the last equality is the MC equation for~$\Theta_\cA$.
\end{proof}

The flatness of the shadow connection says that this derived
coinvariant / bar-cohomological package forms a local system on
$\overline{\mathcal{M}}_{g,n}$: parallel transport between any two
points is path-independent (up to homotopy).  For conformal vertex
algebras at integrable affine level, this recovers the usual local
system of conformal blocks
(Tsuchiya--Ueno--Yamada~\cite{TUY89}); more generally, the
shadow-connection construction supplies the flat bar-cohomological
transport package attached to the bar-intrinsic MC element.

\subsection{The KZ connection at genus $0$}
\label{subsec:thqg-intro-kz-genus0}
\index{Knizhnik--Zamolodchikov connection!from shadow}

At genus~$0$, the moduli space $\overline{\mathcal{M}}_{0,n}$ is the
Deligne--Mumford compactification of the space of $n$-pointed rational
curves.  The shadow connection restricted to genus~$0$ takes the
explicit form:

\begin{equation}\label{eq:thqg-intro-kz-explicit}
\nabla^{\mathrm{KZ}}_{0,n}
\;=\;
d \;-\; \sum_{1 \le i < j \le n}
\frac{r_{ij}(z_i - z_j)}{z_i - z_j}\,dz_i,
\end{equation}
where $r_{ij}(z) = (\operatorname{id} \otimes \operatorname{id}) \circ
r(z) \in \operatorname{End}(M_i) \otimes \operatorname{End}(M_j)$ is
the $r$-matrix acting on the $i$-th and $j$-th module insertions.

When the $r$-matrix is the rational $r$-matrix of a simple Lie
algebra~$\mathfrak{g}$ (that is, $r(z) = \Omega/z$ where
$\Omega = \sum_a t^a \otimes t_a$ is the Casimir element), the
connection~\eqref{eq:thqg-intro-kz-explicit} reduces to the
classical KZ connection:
\begin{equation}\label{eq:thqg-intro-kz-classical}
\nabla^{\mathrm{KZ}}
\;=\;
d \;-\; \frac{1}{k + h^\vee}
\sum_{i < j} \frac{\Omega_{ij}}{z_i - z_j}\,dz_i.
\end{equation}
The level-dependent prefactor $(k + h^\vee)^{-1}$ arises from
the Sugawara normalization of the affine OPE:
$r(z) = \Omega / ((k + h^\vee)\,z)$ for
$\cA = \widehat{\mathfrak{g}}_k$.  This is \emph{not} the
inverse of the modular characteristic
$\kappa(\widehat{\mathfrak{g}}_k)
= \dim\mathfrak{g}\cdot(k + h^\vee)/(2h^\vee)$, which differs
by the factor $\dim\mathfrak{g}/(2h^\vee)$.

The flatness of the KZ connection, originally proved by Knizhnik and
Zamolodchikov~\cite{KZ84} using current algebra Ward identities, is
here a special case of Proposition~\ref{prop:thqg-intro-flatness}: it
follows from the MC equation for~$\Theta_\cA$ restricted to genus~$0$.
On this affine comparison surface, the KZ equations on conformal
blocks are the parallel-transport equations for this flat
connection.

\subsection{The BPZ equations from the shadow obstruction tower}
\label{subsec:thqg-intro-bpz}
\index{BPZ equations!from shadow tower}

For the Virasoro algebra $\cA = \mathrm{Vir}_c$, on the genus-$0$
degenerate-module comparison surface with $n + 1$ marked points
(one carrying a degenerate module $M_{r,s}$ and the rest carrying
generic modules), the shadow connection gives the Virasoro conformal
Ward connection, and the higher collision residues produce the BPZ
differential equations.

The mechanism is as follows.  The degenerate module $M_{r,s}$ has a
null vector at level $rs$; this null vector produces a polynomial
differential operator $\mathcal{D}_{r,s}$ of order~$rs$ in the
insertion variable~$z_0$.  The flatness of the shadow connection
implies that the conformal block
$\Phi(z_0, z_1, \ldots, z_n)$ satisfies
\begin{equation}\label{eq:thqg-intro-bpz-equation}
\mathcal{D}_{r,s}\,\Phi(z_0, z_1, \ldots, z_n)
\;=\; 0.
\end{equation}
At the simplest level $(r,s) = (2,1)$, the null vector is at level~$2$
and the BPZ equation is the second-order PDE
\begin{equation}\label{eq:thqg-intro-bpz-21}
\left[
  \frac{1}{2(2h_{2,1} + 1)}\,\partial_{z_0}^2
  \;-\;
  \sum_{i=1}^n \left(
    \frac{h_i}{(z_0 - z_i)^2}
    + \frac{1}{z_0 - z_i}\,\partial_{z_i}
  \right)
\right]\Phi
\;=\; 0,
\end{equation}
where $h_i$ are the conformal weights of the modules at the
remaining insertions.

From the bar-complex perspective, the BPZ equations arise as follows.
The null vector in~$M_{r,s}$ is a cycle in the bar complex
$\barB(\mathrm{Vir}_c; M_{r,s})$ of the Virasoro with module
coefficients.  The shadow connection transports this cycle across the
moduli space.  The flatness condition
$(\nabla^{\mathrm{hol}})^2 = 0$ then implies that the transport is
consistent, and the transported cycle yields the corresponding BPZ
null-vector equation at each
point of~$\overline{\mathcal{M}}_{0,n+1}$.

The entire BPZ programme (null vectors, differential equations, fusion
rules, minimal model classification) is thus a shadow of the MC
equation in the Virasoro sector of the modular convolution algebra.

\begin{remark}[The BPZ-shadow dictionary]
\label{rem:thqg-intro-bpz-shadow-dictionary}
\index{BPZ equations!shadow dictionary}
The dictionary between classical BPZ data and shadow-tower data is:
\begin{center}
\small
\begin{tabular}{ll}
\textbf{BPZ datum} & \textbf{Shadow datum} \\ \hline
Null vector in $M_{r,s}$ at level $rs$
  & bar cycle in $\barB(\mathrm{Vir}_c; M_{r,s})$ \\[2pt]
BPZ differential operator $\mathcal{D}_{r,s}$
  & shadow connection restricted to degenerate insertion \\[2pt]
Fusion rules $N_{ij}{}^k$
  & dimensions of spaces of conformal blocks \\[2pt]
Minimal model central charges $c_{p,q}$
  & poles of shadow invariants in the $c$-plane \\[2pt]
Kac determinant vanishing locus
  & singular locus of the shadow connection
\end{tabular}
\end{center}
\end{remark}

\begin{remark}[Regularity of the shadow connection]
\label{rem:thqg-intro-shadow-regularity}
\index{shadow connection!regularity}
The shadow connection $\nabla^{\mathrm{hol}}_{g,n}$ is regular singular
along the boundary divisors of $\overline{\mathcal{M}}_{g,n}$: the
connection form has at most simple poles along the boundary strata
where marked points collide or handles degenerate.  This regularity
follows from the structure of the MC element: the amplitude
$\Theta_{\cA,\Gamma}$ on a stable graph~$\Gamma$ has poles of order
at most~$1$ along each boundary divisor, because the bar differential
involves only residue extraction (which produces at most simple poles
from the logarithmic forms~$\eta_{ij}$).

The regular-singular property guarantees that the parallel transport
of the derived coinvariant / protected-state package extends across
the boundary of $\overline{\mathcal{M}}_{g,n}$.  On benchmark
comparison surfaces, this becomes the corresponding extension of
parallel transport for conformal blocks.  The monodromy around a boundary
divisor (a degeneration of the underlying curve) is controlled by the
local exponent of the connection, which is determined by the
collision residue of the $r$-matrix.
\end{remark}

\subsection{Higher-genus extensions}
\label{subsec:thqg-intro-higher-genus-blocks}
\index{derived coinvariants!higher genus}
\index{protected states!higher genus}

At genus $g \ge 1$, the shadow connection acquires additional terms
from the period structure of the curve.  The genus-$1$ shadow
connection is
\begin{equation}\label{eq:thqg-intro-genus1-shadow}
\nabla^{\mathrm{hol}}_{1,n}
\;=\;
\nabla^{\mathrm{hol}}_{0,n}
\;-\;
\kappa(\cA) \cdot \omega_1
\;-\;
\sum_{r \ge 3}
\operatorname{Sh}_{1,n}^{(r)}(\Theta_\cA),
\end{equation}
where $\omega_1$ is the Arakelov form on $\overline{\mathcal{M}}_{1,n}$,
$\kappa(\cA)$ is the modular characteristic, and the higher-arity
shadows $\operatorname{Sh}_{1,n}^{(r)}$ are the contributions from
the shadow obstruction tower at arity~$r$ and genus~$1$.

The term $\kappa(\cA) \cdot \omega_1$ is the curvature of the Hodge
bundle, coupling the derived coinvariant / protected-state package to
the complex structure of the elliptic curve.  On benchmark comparison
surfaces, this is the corresponding coupling of conformal blocks.
The higher-arity corrections encode the nonlinear
dependence of this package on the modular parameter~$\tau$; on
comparison surfaces, of conformal blocks: the
cubic shadow controls the first subleading correction, the quartic
shadow controls the next, and so on.

For the Heisenberg algebra, only the $\kappa$ term survives (the shadow
tower terminates at arity~$2$), and on the standard free-boson
comparison surface the genus-$1$ conformal blocks transform by
multiplication by $\eta(\tau)^{-k}$, where $\eta$ is the
Dedekind eta function.  For the Virasoro algebra, on the torus
comparison surface the infinite shadow obstruction tower produces the modular
differential equation satisfied by Virasoro conformal blocks.

At genus~$g \ge 2$, the shadow connection acquires contributions from
all stable graphs of type $(g,n)$.  The number of contributing graphs
grows with~$g$; the graph sum is controlled by the shadow algebra
$\cA^{\mathrm{sh}}$ (Definition~\ref{def:shadow-algebra}), and the
convergence of the sum is guaranteed by the all-arity master equation
\eqref{eq:all-arity-master-eq-intro}.  The genus-$2$ shadow connection
requires graphs with up to one internal edge (genus-$2$ irreducible
graphs and products of genus-$1$ self-loops), and the genus-$3$
connection requires graphs with up to two internal edges.  The
structure of these contributions is governed by the genus spectral
sequence: the $E_1$ page at genus~$g$ isolates the tree-level ($p = 0$),
one-loop ($p = 1$), and higher-loop ($p \ge 2$) contributions,
with differentials that are the obstruction maps.


% ======================================================================
%  SECTION 4: GRAVITATIONAL COMPLEXITY AND THE EFT HIERARCHY
% ======================================================================

\section{Gravitational complexity and the EFT hierarchy}
\label{sec:thqg-intro-gravitational-complexity}
\index{shadow depth!as gravitational complexity}
\index{EFT hierarchy|textbf}
\index{obstruction classes!and EFT}

The shadow obstruction tower
$\Theta_\cA^{\le 2} \to \Theta_\cA^{\le 3} \to \Theta_\cA^{\le 4}
\to \cdots$
assigns to each chiral algebra a numerical invariant: the
\emph{shadow depth} $r_{\max}(\cA)$, the smallest arity at which the
tower terminates (or $\infty$ if it does not).  This invariant measures
the \emph{gravitational complexity} of~$\cA$: the number of independent
multi-particle interactions required to determine the full modular
structure.

\subsection{Shadow-depth classes}
\label{subsec:thqg-intro-four-classes}
\index{shadow depth!classes}
\index{Gaussian archetype|see{shadow depth}}
\index{Lie/tree archetype|see{shadow depth}}
\index{contact archetype|see{shadow depth}}
\index{mixed archetype|see{shadow depth}}

In general, the shadow depth $r_{\max}(\cA)$ can take any value in
$\{2,3,4,5,\ldots\}\cup\{\infty\}$: a chiral algebra of class
$\mathbf{F}_r$ has $r_{\max} = r$ (tower terminates at arity~$r$),
while class~$\mathbf{M}$ has $r_{\max} = \infty$ (tower does not
terminate).
There is no structural obstruction to any finite value; the
definition is
$r_{\max}=\sup\{r\ge 2:\cA^{\mathrm{sh}}_{r,0}\ne 0\}$.

Within the standard landscape, only the values
$r_{\max}\in\{2,3,4,\infty\}$ are realised
(Definition~\ref{def:shadow-depth-classification}):

\begin{center}
\small
\begin{tabular}{llll}
\textbf{Class} & \textbf{$r_{\max}$} & \textbf{Archetype}
  & \textbf{Examples} \\ \hline
$\mathbf{F}_2=$ \textbf{G} (Gaussian)
  & $2$ & abelian, no bracket
  & Heisenberg, free fermion, lattice \\[2pt]
$\mathbf{F}_3=$ \textbf{L} (Lie/tree)
  & $3$ & Lie bracket, tree-level
  & affine Kac--Moody $\widehat{\mathfrak{g}}_k$ \\[2pt]
$\mathbf{F}_4=$ \textbf{C} (contact/quartic)
  & $4$ & quartic contact interaction
  & $\beta\gamma$, symplectic fermion \\[2pt]
\textbf{M} (mixed)
  & $\infty$ & infinite tower
  & Virasoro, $\mathcal{W}_N$ ($N \ge 3$)
\end{tabular}
\end{center}
The gap between $\mathbf{F}_4$ and $\mathbf{M}$ is an empirical
feature of these families; classes
$\mathbf{F}_5,\mathbf{F}_6,\ldots$ are structurally permitted and
may be realised by chiral algebras outside the standard landscape
(e.g.\ certain non-principal W-algebras or non-freely-generated
vertex algebras).

The classes are determined by the vanishing pattern of the
obstruction sequence:
\begin{alignat}{3}
\textbf{G}: &\quad o_3 = 0, &\quad o_4 &= 0, &\quad o_5 &= 0, \quad \ldots
\label{eq:thqg-intro-vanishing-G} \\
\textbf{L}: &\quad o_3 \ne 0, &\quad o_4 &= 0, &\quad o_5 &= 0, \quad \ldots
\label{eq:thqg-intro-vanishing-L} \\
\textbf{C}: &\quad o_3 \ne 0, &\quad o_4 &\ne 0, &\quad o_5 &= 0, \quad \ldots
\label{eq:thqg-intro-vanishing-C} \\
\mathbf{F}_r: &\quad o_r \ne 0, &\quad o_{r+1} &= 0, &\quad o_{r+2} &= 0, \quad \ldots
\quad\text{(general finite)}
\label{eq:thqg-intro-vanishing-Fr} \\
\textbf{M}: &\quad o_3 \ne 0, &\quad o_4 &\ne 0, &\quad o_5 &\ne 0, \quad \ldots
\label{eq:thqg-intro-vanishing-M}
\end{alignat}
(In class~\textbf{L}, $o_3 \ne 0$ means the cubic shadow is present
but the quartic obstruction vanishes; the tower terminates at
arity~$3$.  In class~\textbf{C}, the quartic is the last nonzero
obstruction.  In general class~$\mathbf{F}_r$, the tower terminates
at arity~$r$.  In class~\textbf{M}, all obstructions are nonzero.)

\subsection{The obstruction classes as effective field theory data}
\label{subsec:thqg-intro-eft-interpretation}
\index{effective field theory!from shadow tower}

Each obstruction class $o_{r+1}(\cA)$ has a physical interpretation as
an effective interaction vertex.  The shadow obstruction tower is an
effective field theory (EFT) hierarchy:

\begin{enumerate}[label=\textup{(\roman*)}]
\item \emph{Arity $2$: the free-field kernel.}
  The modular characteristic $\kappa(\cA)$ is the one-loop determinant
  of the free theory.  It counts the number of degrees of freedom
  weighted by spin: $\kappa(\cH_k) = k$ for bosons,
  $\kappa(\psi_k) = k/4$ for $k$ free fermions
  ($c = k/2$, so $\kappa = c/2 = k/4$).  In the EFT language,
  $\kappa$ is the cosmological constant of the two-dimensional
  gravitational theory.

\item \emph{Arity $3$: the cubic interaction.}
  The cubic shadow $\mathfrak{C}(\cA)$ is the three-point coupling.
  For affine Kac--Moody, it is proportional to the structure constants
  $f^{abc}$ of the Lie algebra; for the Virasoro, it is proportional to
  the central charge.  The cubic shadow vanishes if and only if the
  OPE has no simple poles; the algebra is abelian.

\item \emph{Arity $4$: the quartic resonance.}
  The quartic class $\mathfrak{Q}(\cA)$ is the four-point contact
  interaction.  For the Virasoro algebra, the quartic contact
  invariant is
  \[
  Q^{\mathrm{contact}}_{\mathrm{Vir}}
  \;=\; \frac{10}{c(5c+22)},
  \]
  with poles at $c = 0$ and $c = -22/5$ (the $(2,5)$ minimal model
  central charge).  The quartic resonance class carries a clutching
  law under stable-curve degeneration
  (Theorem~\ref{thm:nms-clutching-law-modular-resonance}).

\item \emph{Arity $5$ and beyond: the quintic forced.}
  For the Virasoro and $\mathcal{W}_N$ algebras, $o_5(\cA) \ne 0$:
  the quintic interaction is forced, and the tower never terminates.
  Each higher arity~$r$ introduces a new effective vertex, coupled to
  the lower-order vertices by the all-arity master equation
  $\nabla_H(\mathrm{Sh}_r) + o^{(r)} = 0$.
\end{enumerate}

The EFT hierarchy is organized by a single principle: the shadow
obstruction tower is the Taylor expansion of the full MC element
$\Theta_\cA$ in powers of the arity variable.  The radius of
convergence is the shadow depth~$r_{\max}$.  For Gaussian algebras,
the expansion converges at linear order; for Lie/tree algebras, at
cubic order; for contact algebras, at quartic order; for mixed
algebras, the expansion has infinite radius and the full MC element
cannot be captured by any finite truncation.

\subsection{The genus loop and primitive shadows}
\label{subsec:thqg-intro-genus-loop}
\index{genus loop!and primitive shadows}

The descent from arity~$r$ to arity~$r-2$ is governed by the
\emph{genus loop}
\begin{equation}\label{eq:thqg-intro-genus-loop}
\Lambda_P \colon \operatorname{Sym}^r
\;\longrightarrow\;
\operatorname{Sym}^{r-2},
\end{equation}
which contracts a pair of indices using the self-loop graph
(the genus-$1$ sewing operation at a single vertex).
The genus loop reduces the arity by~$2$: from the quartic shadow,
it produces a scalar correction to the quadratic level (a correction
to~$\kappa$); from the cubic shadow, it produces a correction to
the vacuum (a scalar).

The kernel of the genus loop consists of the \emph{primitive shadows}:
the irreducible pieces that cannot be obtained by trace contraction
from higher arity.  The primitive-shadow decomposition
\begin{equation}\label{eq:thqg-intro-primitive-decomposition}
\operatorname{Sh}_r(\Theta_\cA)
\;=\;
\operatorname{Sh}_r^{\mathrm{prim}}(\Theta_\cA)
\;+\;
\Lambda_P\!\left(
  \operatorname{Sh}_{r+2}^{\mathrm{prim}}(\Theta_\cA)
\right)
\;+\;
\Lambda_P^2\!\left(
  \operatorname{Sh}_{r+4}^{\mathrm{prim}}(\Theta_\cA)
\right)
\;+\; \cdots
\end{equation}
expresses the full shadow at arity~$r$ as a sum of primitive
contributions from higher arities, contracted downward by the
genus loop.  This is the modular analogue of the
primitive-element decomposition in a Hopf algebra.

\subsection{The operadic complexity theorem}
\label{subsec:thqg-intro-operadic-complexity}
\index{operadic complexity theorem}
\index{shadow depth!and $A_\infty$-depth}

The shadow depth equals the $A_\infty$-depth of the bar cohomology:

\begin{theorem}[Operadic complexity; \ClaimStatusProvedHere;
Theorem~\ref{thm:operadic-complexity}]
\label{thm:thqg-intro-operadic-complexity}% NOTE: retains conj: prefix for backward compatibility; claim is ProvedHere
For any chiral algebra~$\cA$ on~$X$,
\[
r_{\max}(\cA) \;=\; \text{$A_\infty$-depth of $H^*(\barB_X(\cA))$}
\;=\; \text{$L_\infty$-formality level of $\gAmod$}.
\]
\end{theorem}

The $A_\infty$-depth is the smallest~$r$ such that the transferred
$A_\infty$-product $m_r$ on bar cohomology is nonzero.  The
$L_\infty$-formality level is the smallest~$r$ such that the minimal
$L_\infty$-model of the modular convolution algebra has a nonzero
$r$-ary bracket.  The proof proceeds in two steps: $r_{\max} = f_\infty$
is Theorem~\ref{thm:shadow-formality-identification}
(Proposition~\ref{prop:shadow-massey-identification} gives the
genus-$0$ case; the all-genera direction uses
Proposition~\ref{prop:shadow-formality-low-arity}, proved at
arities~$2$, $3$, and~$4$, extended by the HPL tree formula to all
finite arities); $f_\infty = d_\infty$ follows because antisymmetrization
is injective on cyclic cochains.  The full argument is
Theorem~\ref{thm:operadic-complexity-detailed}.

\subsection{Shadow depth versus Koszulness}
\label{subsec:thqg-intro-depth-vs-koszulness}
\index{shadow depth!versus Koszulness}
\index{Koszulness!not determined by shadow depth}

Shadow depth does \emph{not} characterize Koszulness.  Both finite and
infinite shadow-depth algebras can be Koszul:
\begin{itemize}
\item Heisenberg ($r_{\max} = 2$): Koszul.
\item Affine Kac--Moody ($r_{\max} = 3$): Koszul.
\item $\beta\gamma$ ($r_{\max} = 4$): Koszul.
\item Virasoro ($r_{\max} = \infty$): Koszul.
\end{itemize}
Shadow depth classifies the \emph{complexity} of Koszul algebras, not
Koszulness itself.  The Koszulness characterization programme
(Theorem~\ref{thm:koszul-equivalences-meta}: ten unconditional
equivalences, one conditional, one one-directional) is independent
of the shadow depth classification.

The distinction is precise: Koszulness means the bar spectral sequence
degenerates at~$E_2$; shadow depth measures how many multi-particle
operations the bar differential requires.  The Virasoro is Koszul
(its bar cohomology is concentrated in bar degree~$1$) but has
infinite shadow depth (its $A_\infty$ minimal model has nonzero
products at all arities).  These are compatible because Koszulness
concerns the cohomological concentration of the bar complex, while
shadow depth concerns the complexity of the bar differential.


% ======================================================================
%  SECTION 5: DUALITY AND THE GRAVITATIONAL PHASE SPACE
% ======================================================================

\section{Duality and the gravitational phase space}
\label{sec:thqg-intro-duality-phase-space}
\index{Verdier duality!and gravitational phase space}
\index{complementarity!shifted-symplectic}
\index{Lagrangian!complementary pair}
\index{critical locus}

Koszul duality on a curve is not a symmetry of the chiral algebra but
a symmetry of the genus tower.  The Verdier involution
$\sigma = \mathbb{D}_{\operatorname{Ran}}$ acts on the ambient complex
$\mathbf{C}_g(\cA)$ and decomposes it into complementary halves.  The
resulting structure is a shifted-symplectic phase space whose
Lagrangian polarization records the gravitational content of Koszul
duality.

\subsection{The Verdier involution and the ambient decomposition}
\label{subsec:thqg-intro-verdier-decomposition}
\index{Verdier duality!involution}

The Verdier duality functor
$\mathbb{D}_{\operatorname{Ran}} \colon
D(\mathcal{D}\text{-mod}(\operatorname{Ran}(X)))
\to D(\mathcal{D}\text{-mod}(\operatorname{Ran}(X)))$
sends the bar coalgebra to the homotopy Koszul dual algebra
$\mathbb{D}_{\operatorname{Ran}}\,\barB_X(\cA) \simeq \cA^!_\infty$
(Theorem~A; factorization \emph{algebra}, not coalgebra).
The induced involution on the ambient complex is
$\sigma \colon \mathbf{C}_g(\cA) \to \mathbf{C}_g(\cA)$, and the
eigenspace decomposition is
\begin{equation}\label{eq:thqg-intro-verdier-eigenspaces}
\mathbf{C}_g(\cA)
\;\simeq\;
\mathbf{Q}_g(\cA) \;\oplus\; \mathbf{Q}_g(\cA^!),
\end{equation}
where $\mathbf{Q}_g(\cA) = \operatorname{fib}(\sigma - \operatorname{id})$
and $\mathbf{Q}_g(\cA^!) = \operatorname{fib}(\sigma + \operatorname{id})$
are the $+1$ and $-1$ eigenspaces.  This is Theorem~C
(Theorem~\ref{thm:quantum-complementarity-main}).

The decomposition is \emph{orthogonal} for the shifted-symplectic
pairing: the two eigenspaces are isotropic, and the pairing between
them is perfect.  The ambient complex is therefore a
\emph{symplectic vector space} (in the shifted-derived sense), and
the Verdier involution is the map that exchanges the two Lagrangian
halves.

\subsection{The shifted-symplectic structure}
\label{subsec:thqg-intro-shifted-symplectic}
\index{shifted-symplectic!structure on ambient complex}

The ambient complex carries a canonical shifted-symplectic structure.
The pairing
\begin{equation}\label{eq:thqg-intro-shifted-symplectic-pairing}
\omega_{g}^{\mathrm{symp}} \colon
\mathbf{C}_g(\cA) \otimes \mathbf{C}_g(\cA)
\;\longrightarrow\;
\mathbb{C}[-(3g-3)]
\end{equation}
is a $(-(3g-3))$-shifted symplectic form on the derived moduli of
genus-$g$ bar families, in the sense of
Pantev--To\"en--Vaqui\'e--Vezzosi~\cite{PTVV13}
\textup{(}Proposition~\ref{prop:ptvv-lagrangian}\textup{)}.  The shift
$-(3g-3)$ reflects the dimension of the moduli space
$\overline{\mathcal{M}}_g$: the symplectic form lives in the top
cohomological degree of~$\overline{\mathcal{M}}_g$.

\subsection{Lagrangian complementarity}
\label{subsec:thqg-intro-lagrangian}
\index{Lagrangian!complementarity}

The eigenspaces $\mathbf{Q}_g(\cA)$ and $\mathbf{Q}_g(\cA^!)$ are
\emph{Lagrangian} for the shifted-symplectic pairing: each is an
isotropic subspace of half the total rank.  The Lagrangian condition
means that
\begin{equation}\label{eq:thqg-intro-lagrangian-condition}
\omega_g^{\mathrm{symp}}\bigl|_{\mathbf{Q}_g(\cA) \otimes \mathbf{Q}_g(\cA)}
\;=\; 0,
\qquad
\omega_g^{\mathrm{symp}}\bigl|_{\mathbf{Q}_g(\cA^!) \otimes \mathbf{Q}_g(\cA^!)}
\;=\; 0,
\end{equation}
and the pairing between the two Lagrangians is a perfect duality:
\begin{equation}\label{eq:thqg-intro-lagrangian-perfect}
\omega_g^{\mathrm{symp}} \colon
\mathbf{Q}_g(\cA) \otimes \mathbf{Q}_g(\cA^!)
\;\xrightarrow{\;\sim\;}
\mathbb{C}[-(3g-3)].
\end{equation}

The Lagrangian pair $(\mathbf{Q}_g(\cA), \mathbf{Q}_g(\cA^!))$ is
the gravitational phase space of the chiral algebra: the two
Lagrangians correspond to the boundary and the bulk degrees of
freedom respectively, and the symplectic pairing is the holographic
coupling between them.  Koszul duality exchanges the two
Lagrangians: $\mathbb{D}(\mathbf{Q}_g(\cA)) \simeq \mathbf{Q}_g(\cA^!)$.

\subsection{The critical locus and anomaly cancellation}
\label{subsec:thqg-intro-critical-locus}
\index{critical locus!and anomaly}
\index{anomaly cancellation!from Koszul duality}

The modular characteristic transforms under Koszul duality as
$\kappa(\cA) + \kappa(\cA^!) = 0$ for affine Kac--Moody and free-field
algebras; for $\mathcal{W}$-algebras the complementarity sum
$\kappa + \kappa'$ is a nonzero constant
(e.g.\ $13$ for Virasoro, $250/3$ for $\mathcal{W}_3$; see
Theorem~\ref{thm:modular-characteristic}).  The \emph{critical locus} is the
locus where $\kappa(\cA) = 0$: the genus-$1$ curvature vanishes, and
the bar complex is nilpotent at all genera simultaneously.

On the critical locus, the shifted-symplectic pairing degenerates: the
Lagrangian decomposition persists but the genus-$1$ obstruction
$\mathrm{obs}_1(\cA) = \kappa(\cA) \cdot \lambda_1$ vanishes, and the
ambient complex splits into a direct sum without curvature correction.
Physically, the critical locus is the anomaly-free locus: the theory
propagates on higher-genus surfaces without obstruction from the Hodge
bundle.

For affine Kac--Moody, the anti-symmetry $\kappa(\widehat{\mathfrak{g}}_k)
+ \kappa(\widehat{\mathfrak{g}}_{k'}) = 0$ with $k' = -k - 2h^\vee$ means
that the critical locus is $k = -h^\vee$: the critical level where the
Sugawara construction is undefined.  For the Virasoro algebra,
$\kappa(\mathrm{Vir}_c) + \kappa(\mathrm{Vir}_{26-c}) = c/2 + (26-c)/2 = 13$ by the
complementarity $c + c' = 26$; the chiral Koszul self-dual point is $c = 13$
($\mathrm{Vir}_{13}^! = \mathrm{Vir}_{13}$, $\kappa = 13/2$).  For the $\beta\gamma$ system
at conformal weight $\lambda$,
$\kappa(\beta\gamma) = c_{\beta\gamma}/2 = 6\lambda^2 - 6\lambda + 1$ and $\kappa(bc) = -\kappa(\beta\gamma)$; the
tensor product
$\beta\gamma \otimes bc$ has $\kappa = 0$.

\subsection{The complementarity sum}
\label{subsec:thqg-intro-complementarity-sum}
\index{complementarity sum}

The Koszul conductor $K(\cA) = c(\cA) + c(\cA^!)$ is the sum of
central charges of the Koszul pair.  The conductor is determined by
the duality constraint on~$\kappa$
(Theorem~\ref{thm:modular-characteristic}(iii)) together with the
Sugawara formula/Drinfeld--Sokolov formula for the central charge.  The
explicit values are:
\begin{center}
\small
\begin{tabular}{lll}
\textbf{Family} & \textbf{Conductor $K = c + c'$}
  & \textbf{Key values} \\ \hline
Heisenberg $\cH_k$ & $2$ & all $k$ \\[2pt]
Affine $\widehat{\mathfrak{g}}_k$ & $2\dim\mathfrak{g}$
  & $6$ for $\mathfrak{sl}_2$ \\[2pt]
Principal $\mathcal{W}^k(\mathfrak{g})$
  & $2r + 4h^\vee d$ & $26$ for $\mathfrak{sl}_2$ \\[2pt]
$\beta\gamma$ & $-2$ & \\[2pt]
Virasoro $\mathrm{Vir}_c$ & $26$ & $c = 13$ Koszul self-dual
\end{tabular}
\end{center}
The conductor $K = 26$ for the Virasoro family recovers the bosonic
string critical dimension
(Theorem~\ref{thm:central-charge-complementarity}).

\subsection{Duality of the shadow obstruction tower}
\label{subsec:thqg-intro-duality-shadow}
\index{shadow obstruction tower!duality}

The Verdier involution acts on the full MC element by
$\mathbb{D}(\Theta_\cA) = \Theta_{\cA^!}$, and hence on the shadow
obstruction tower level by level.  For affine Kac--Moody and free-field
algebras (where the Verdier involution acts as sign reversal on each
shadow component):
\begin{align}
\kappa(\cA) + \kappa(\cA^!) &= 0, \label{eq:thqg-intro-kappa-duality} \\
\mathfrak{C}(\cA) + \mathfrak{C}(\cA^!) &= 0, \label{eq:thqg-intro-cubic-duality} \\
\mathfrak{Q}(\cA) + \mathfrak{Q}(\cA^!) &= 0. \label{eq:thqg-intro-quartic-duality}
\end{align}
For $\mathcal{W}$-algebras arising from Drinfeld--Sokolov reduction,
$\kappa + \kappa'$ is a nonzero constant determined by the root datum
(e.g.\ $13$ for Virasoro); see
Theorem~\ref{thm:modular-characteristic}(iii).
The duality of the shadow obstruction tower means that the gravitational complexity
of~$\cA$ and~$\cA^!$ are the same: $r_{\max}(\cA) = r_{\max}(\cA^!)$.
The gravitational phase space sees both algebras symmetrically.

\subsection{The Lagrangian intersection and the center}
\label{subsec:thqg-intro-lagrangian-intersection}
\index{Lagrangian!intersection}
\index{center!as Lagrangian intersection}

The intersection of the two Lagrangians,
\begin{equation}\label{eq:thqg-intro-lagrangian-intersection}
\mathbf{Q}_g(\cA) \;\cap\; \mathbf{Q}_g(\cA^!)
\;=\;
\ker(\sigma - \operatorname{id}) \;\cap\; \ker(\sigma + \operatorname{id})
\;=\;
\ker(\sigma^2 - \operatorname{id}),
\end{equation}
is nontrivial only when $\sigma^2 \ne \operatorname{id}$; that is,
when the Verdier involution has nontrivial higher-order structure.
For Koszul pairs on the critical locus ($\kappa = 0$), the involution
is strict ($\sigma^2 = \operatorname{id}$) and the intersection is
the center~$Z(\cA)$ of the chiral algebra: the space of chiral-central
sections that are simultaneously $+1$ and $-1$ under Verdier
involution.  Off the critical locus, the intersection is empty
(the two Lagrangians are transverse).


% ======================================================================
%  SECTION 6: FROM COLLISION TO YANG--BAXTER
% ======================================================================

\section{From collision to Yang--Baxter}
\label{sec:thqg-intro-collision-yang-baxter}
\index{collision filtration|textbf}
\index{Yang--Baxter equation!from collision}
\index{r-matrix@$r$-matrix!from collision filtration}
\index{Yangian!from collision filtration}

The spectral $r$-matrix $r(z)$ in the holographic datum is extracted
from the MC element by a filtration indexed by the depth of collision
in the Fulton--MacPherson compactification.  The classical Yang--Baxter
equation is a formal consequence of the Arnold relation at
codimension~$2$.  The Yangian arises as the endomorphism algebra of
this construction.

\subsection{The collision filtration}
\label{subsec:thqg-intro-collision-definition}
\index{collision filtration!definition}

Let $\overline{C}_n(X)$ be the Fulton--MacPherson compactification of
the configuration space of $n$ distinct points on~$X$.  The boundary
$\partial \overline{C}_n(X)$ is a normal crossing divisor whose
components are indexed by subsets $S \subset \{1, \ldots, n\}$ with
$|S| \ge 2$: the stratum $D_S$ parametrizes configurations where the
points in~$S$ collide.

Define the \emph{collision depth} of a boundary stratum as the number
of nested collisions:
\begin{equation}\label{eq:thqg-intro-collision-depth}
\operatorname{depth}(D_{S_1 \subset \cdots \subset S_k})
\;=\; k.
\end{equation}
The collision filtration on the bar complex is
\begin{equation}\label{eq:thqg-intro-collision-filtration}
F^0 \barB_X(\cA)
\;\supset\;
F^1 \barB_X(\cA)
\;\supset\;
F^2 \barB_X(\cA)
\;\supset\; \cdots,
\end{equation}
where $F^p \barB_X(\cA)$ consists of bar elements supported on strata
of collision depth $\ge p$.  The associated graded is
\[
\operatorname{gr}^p_F \barB_X(\cA)
\;\cong\;
\bigoplus_{|S| = p+1}
\barB_{T_{z_0}X}(\cA|_S) \otimes \barB_X(\cA|_{S^c \cup \{z_0\}}),
\]
where the first factor records the infinitesimal collision pattern
(the ``screen'') and the second factor records the residual
configuration.

The collision filtration is multiplicative: the bar coproduct respects
the filtration, and the associated graded inherits a coalgebra structure.
The differential on the associated graded is the tree-level
approximation to the bar differential: it captures the leading-order
OPE at each collision depth.

\subsection{The spectral parameter}
\label{subsec:thqg-intro-spectral-parameter}
\index{spectral parameter!from collision}

At collision depth~$1$, two points collide: $z_1 \to z_2$.  The
screen is $T_{z_0} X \cong \mathbb{C}$, parametrized by the relative
coordinate $z = z_1 - z_2$.  The bar differential restricted to
depth~$1$ produces the binary OPE, expressed as a Laurent series
in~$z$:
\begin{equation}\label{eq:thqg-intro-spectral-r}
r(z)
\;=\;
\operatorname{gr}^1_F(d_{\barB})|_{\barB^2}
\;=\;
\sum_{n \ge 0} \frac{r_n}{z^{n+1}},
\end{equation}
where $r_n \in \cA^! \otimes \cA^!$ is the $n$-th Laurent coefficient,
encoding the $n$-th OPE mode.  The element
$r(z) \in \cA^! \,\widehat{\otimes}\, \cA^!(\!(z^{-1})\!)$
is the \emph{spectral $r$-matrix}: a meromorphic endomorphism of
the tensor square of the Koszul dual, with a pole at $z = 0$
(the collision locus) and regular as $z \to \infty$ (the separation
limit).

The residue of $r(z)$ at $z = 0$ is the leading OPE mode:
$\operatorname{Res}_{z=0}\, r(z) = r_0 \in \cA^! \otimes \cA^!$.
For affine Kac--Moody, $r_0 = \Omega = \sum_a t^a \otimes t_a$ is
the Casimir element.  The higher Laurent coefficients $r_1, r_2, \ldots$
encode the subleading OPE modes and determine the higher-spin structure
of the dg-shifted Yangian.

\subsection{Arnold implies CYBE}
\label{subsec:thqg-intro-arnold-cybe}
\index{Arnold relation!implies CYBE}
\index{classical Yang--Baxter equation}

The classical Yang--Baxter equation (CYBE) for~$r(z)$ is a
formal consequence of the Arnold relation on~$\overline{C}_3(X)$.
The argument proceeds as follows.

At bar degree~$3$, the MC equation $D\Theta + \frac{1}{2}[\Theta,\Theta] = 0$
restricted to genus~$0$ and arity~$3$ gives
\begin{equation}\label{eq:thqg-intro-mc-arity3}
d_{\barB}(\Theta^{(3)})
+ \tfrac{1}{2}[\Theta^{(2)}, \Theta^{(2)}]^{(3)}
\;=\; 0,
\end{equation}
where $\Theta^{(r)}$ denotes the arity-$r$ component.  The second
term is the operadic composition of two binary operations, which at
collision depth~$2$ produces the triple collision residue.

The collision filtration at depth~$2$ restricts to the codimension-$2$
stratum $D_{ijk}$ of triple collision in $\overline{C}_3(X)$.  The
three pairs $(ij)$, $(jk)$, $(ik)$ produce three terms in the
iterated residue:
\begin{align}
&[r_{12}(z_1 - z_2),\, r_{13}(z_1 - z_3)] \nonumber \\
+\; &[r_{12}(z_1 - z_2),\, r_{23}(z_2 - z_3)] \label{eq:thqg-intro-cybe-explicit} \\
+\; &[r_{13}(z_1 - z_3),\, r_{23}(z_2 - z_3)]
\;=\; 0. \nonumber
\end{align}
Each term corresponds to a boundary stratum of $\overline{C}_3(X)$,
and the cyclic sum vanishes because of the Arnold relation
$\eta_{12} \wedge \eta_{23} + \eta_{23} \wedge \eta_{31}
+ \eta_{31} \wedge \eta_{12} = 0$ at the codimension-$2$ intersection.
The vanishing of this cyclic sum is exactly the CYBE for~$r(z)$.

The derivation is Theorem~\ref{thm:collision-depth-2-ybe}.  The key
point is that no input beyond the geometry of $\overline{C}_3(X)$ is
needed: the CYBE is a topological identity, inherited from the
Arnold relation on configuration spaces.

\subsection{From $r$-matrix to Yangian}
\label{subsec:thqg-intro-r-to-yangian}
\index{Yangian!construction from $r$-matrix}
\index{dg-shifted Yangian}

The $r$-matrix $r(z)$ determines a braided tensor structure on the
line-operator category $\mathcal{C}(\cA) = \cA^!\text{-}\mathsf{mod}$.
The endomorphism algebra of this tensor structure is the Yangian.

\begin{definition}[The dg-shifted Yangian]
\label{def:thqg-intro-dg-yangian}
\index{dg-shifted Yangian!definition}
The \emph{dg-shifted Yangian} $\mathcal{Y}^{\mathrm{dg}}_T(\cA)$
is the dg algebra generated by the modes of the $r$-matrix, with
relations imposed by the CYBE and the higher collision conditions:
\begin{equation}\label{eq:thqg-intro-yangian-generators}
\mathcal{Y}^{\mathrm{dg}}_T(\cA)
\;=\;
\operatorname{End}_{\otimes_{r(z)}}\!
\bigl(\mathcal{C}(\cA)\bigr),
\end{equation}
the endomorphism algebra of the tensor functor $\otimes_{r(z)}$ on
$\cA^!\text{-}\mathsf{mod}$.
\end{definition}

For affine Kac--Moody $\cA = \widehat{\mathfrak{g}}_k$, the
dg-shifted Yangian recovers the Yangian $Y(\mathfrak{g})$: the
rational $r$-matrix $r(z) = \Omega/z$ generates the RTT presentation
of $Y(\mathfrak{g})$ (Theorem~\ref{thm:yangian-bar-rtt}), and the
CYBE becomes the defining relation of the Yangian.  The evaluation
homomorphism $\operatorname{ev}_a \colon Y(\mathfrak{g}) \to
U(\mathfrak{g})$ corresponds to the restriction of the $r$-matrix
to the depth-$1$ collision stratum.

For the Virasoro algebra, the dg-shifted Yangian is the deformation of
$Y(\mathfrak{sl}_2)$ by the Virasoro central charge: the $r$-matrix
has a double pole (from the stress tensor OPE), and the Yangian
generators acquire a quadratic correction proportional to~$c$.

\subsection{The RTT presentation}
\label{subsec:thqg-intro-rtt}
\index{RTT presentation}

The Yangian admits a presentation through the RTT equation.  Define the
\emph{monodromy matrix}
\begin{equation}\label{eq:thqg-intro-monodromy}
T(z) \;=\; \sum_{n \ge 0} \frac{T_n}{z^{n+1}}
\;\in\;
\operatorname{End}(V) \otimes \mathcal{Y}^{\mathrm{dg}}_T(\cA)[\![z^{-1}]\!],
\end{equation}
where $V$ is the evaluation module and $T_n$ are the Yangian generators.
The RTT equation is
\begin{equation}\label{eq:thqg-intro-rtt-equation}
R(z_1 - z_2)\,T_1(z_1)\,T_2(z_2)
\;=\;
T_2(z_2)\,T_1(z_1)\,R(z_1 - z_2),
\end{equation}
where $R(z) = 1 + r(z)/z$ is the quantum $R$-matrix (the classical
$r$-matrix plus the identity), $T_i(z) = T(z)$ acting in the $i$-th
tensor factor, and the equation holds in
$\operatorname{End}(V^{\otimes 2}) \otimes \mathcal{Y}[\![z_1^{-1}, z_2^{-1}]\!]$.

The RTT equation is the quantization of the CYBE: the classical limit
$\hbar \to 0$ of $R = 1 + \hbar r + O(\hbar^2)$ gives the CYBE
for~$r(z)$.  From the bar-complex perspective, the RTT equation arises
from the collision filtration at depth~$3$: the depth-$3$ strata of
$\overline{C}_4(X)$ impose relations among the depth-$2$ generators,
and these relations are exactly the RTT equation.

\subsection{The modular Yangian}
\label{subsec:thqg-intro-modular-yangian}
\index{modular Yangian}

The all-genera extension of the Yangian construction is the
\emph{modular dg-shifted Yangian}
$\mathcal{Y}^{\mathrm{dg,mod}}_T(\cA)$
(Definition~\ref{def:modular-yangian-pro}): the pronilpotent
completion of the dg-shifted Yangian intended to organize genus
corrections.  A conjectural modular $R$-matrix
\begin{equation}\label{eq:thqg-intro-modular-r}
R^{\mathrm{mod}}_T(z; \hbar)
\;=\;
1 + \sum_{g \ge 0} \hbar^{2g} R^{(g)}_T(z)
\end{equation}
would be a Maurer--Cartan element in the completed convolution algebra,
and the corresponding higher-genus Yang--Baxter condition would be the
modular-kernel shadow of the MC equation
$D\Theta + \frac{1}{2}[\Theta,\Theta] = 0$ at the level of the modular
convolution algebra.

At genus~$0$, the modular $R$-matrix reduces to the classical
$r$-matrix: $R^{(0)}_T(z) = r(z)$.  At genus~$1$, the correction
$R^{(1)}_T(z)$ is intended to encode the elliptic deformation: the
$r$-matrix would acquire a dependence on the modular parameter $\tau$
of the torus, controlled by the modular characteristic~$\kappa(\cA)$.
The genus expansion of the modular $R$-matrix is thus only a
conjectural higher-genus refinement of the perturbative holographic
scattering package on this surface.

\subsection{The Drinfeld--Kohno bridge}
\label{subsec:thqg-intro-drinfeld-kohno}
\index{Drinfeld--Kohno theorem!from bar complex}

The Drinfeld--Kohno theorem connects the KZ connection (a flat
connection on the configuration space, with monodromy in the braid
group) to the quantum group (a Hopf algebra whose representation
category is braided monoidal).  From the bar-complex perspective,
this bridge is the passage from the shadow connection
(\S\ref{sec:thqg-intro-shadow-connection}) to the Yangian
(\S\ref{subsec:thqg-intro-r-to-yangian}):
\begin{equation}\label{eq:thqg-intro-dk-bridge}
\begin{tikzcd}[column sep=huge]
\text{KZ connection } \nabla^{\mathrm{KZ}}
  \arrow[r, "\text{monodromy}"]
& \text{Braid-group representation}
  \arrow[r, "\text{DK}"]
& U_q(\mathfrak{g})\text{-}\mathsf{mod}
\end{tikzcd}
\end{equation}
The first arrow computes the monodromy of the KZ connection around
loops in the configuration space; the second arrow identifies this
monodromy representation with the braided tensor category of
$U_q(\mathfrak{g})$-modules.  On the affine Kac--Moody comparison
surface, the first arrow is classified by the bar complex because the
KZ connection is identified there with the genus-$0$ shadow
connection; the second arrow is the Drinfeld--Kohno comparison with the
quantum-group side.

The parameter matching is
$q = e^{\pi i / (k + h^\vee)}$, where $k$ is the level of the affine
Kac--Moody algebra and $h^\vee$ is the dual Coxeter number.  This is
the Kazhdan--Lusztig equivalence
(Theorem~\ref{thm:kazhdan-lusztig-equivalence}): the module category
of $\widehat{\mathfrak{g}}_k$ at non-critical level is equivalent to
the representation category of $U_q(\mathfrak{g})$, with $q$ determined
by the level.


% ======================================================================
%  SECTION 7: PERTURBATIVE FINITENESS AND THE CRITICAL STRING
% ======================================================================

\section{Perturbative finiteness and the critical string}
\label{sec:thqg-intro-perturbative-finiteness}
\index{perturbative finiteness|textbf}
\index{HS-sewing condition|textbf}
\index{A-hat genus@$\hat{A}$-genus!transgression}
\index{critical string!$c = 26$}

The shadow obstruction tower is an algebraic construction: it lives in
the cyclic deformation complex and requires no analysis.  The passage
to physics (partition functions, amplitudes, modular forms) requires
an analytic input: the convergence of the genus tower.  This section
describes the analytic completion programme and its relation to the
critical string.

\subsection{The HS-sewing criterion}
\label{subsec:thqg-intro-hs-sewing}
\index{HS-sewing condition!criterion}
\index{Hilbert--Schmidt!sewing condition}

The analytic convergence of the genus tower is controlled by a
Hilbert--Schmidt condition on the sewing amplitudes.  Let
$\cA$ be a conformal vertex algebra with vacuum character
$\chi_\cA(q) = \operatorname{tr}_\cA(q^{L_0 - c/24})$.  The
\emph{sewing amplitude} at genus~$g$ is the functional
\begin{equation}\label{eq:thqg-intro-sewing-amplitude}
Z_g(\cA; q_1, \ldots, q_g)
\;=\;
\operatorname{tr}_{\cA^{\otimes g}}\!\left(
  \prod_{\alpha=1}^g q_\alpha^{L_0 - c/24}
  \cdot S_\alpha
\right),
\end{equation}
where $S_\alpha$ is the sewing operator for the $\alpha$-th handle.

\begin{definition}[HS-sewing condition]
\label{def:thqg-intro-hs-sewing}
\index{HS-sewing condition!definition}
The chiral algebra~$\cA$ satisfies the \emph{HS-sewing condition} if,
for all $0 < |q| < 1$ and all sewing data, the sum
\begin{equation}\label{eq:thqg-intro-hs-condition}
\sum_{a,b,c} |q|^{a+b+c}\, \|m^c_{a,b}\|_{\mathrm{HS}}^2
\;<\; \infty,
\end{equation}
where $m^c_{a,b}$ is the OPE structure constant between states of
conformal weight $a$, $b$, $c$, and $\|\cdot\|_{\mathrm{HS}}$ is
the Hilbert--Schmidt norm on the finite-dimensional weight spaces.
\end{definition}

The HS-sewing condition implies that the sewing
amplitude~\eqref{eq:thqg-intro-sewing-amplitude} converges absolutely
and defines a holomorphic function of $(q_1, \ldots, q_g)$ in the
polydisk $|q_\alpha| < 1$.  The convergent sum is the genus-$g$
partition function $Z_g(\cA)$.

\begin{theorem}[General HS-sewing criterion;
\ClaimStatusProvedElsewhere]
\label{thm:thqg-intro-hs-general}
Polynomial OPE growth and subexponential sector growth imply the
HS-sewing condition at all genera.  In particular, the HS-sewing
condition holds for every algebra in the standard landscape:
Heisenberg, free fermion, lattice, $\beta\gamma$, affine Kac--Moody,
Virasoro, $\mathcal{W}$-algebras.  The proof is
Theorem~\textup{\ref{thm:general-hs-sewing}}.
\end{theorem}

\subsection{The Heisenberg sewing theorem}
\label{subsec:thqg-intro-heisenberg-sewing}
\index{Heisenberg algebra!sewing theorem}
\index{Fredholm determinant!Heisenberg sewing}

For the Heisenberg algebra, the sewing reduces to a one-particle
computation.  The genus-$g$ partition function is a Fredholm
determinant of the Bergman sewing operator:

\begin{theorem}[Heisenberg sewing;
\ClaimStatusProvedElsewhere]
\label{thm:thqg-intro-heisenberg-sewing}
The genus-$g$ partition function of $\cH_k$ is
\begin{equation}\label{eq:thqg-intro-heisenberg-fredholm}
Z_g(\cH_k)
\;=\;
\det\!\left(1 - S^{(1)}_g\right)^{-k},
\end{equation}
where $S^{(1)}_g$ is the one-particle Bergman sewing operator,
a trace-class operator on the one-particle Hilbert space
$\cH_k^{(1)} = \bigoplus_{n \ge 1} \mathbb{C}\cdot\alpha_{-n}|0\rangle$.
The eigenvalues of~$S^{(1)}_g$ are the squares of the Schottky
parameters.  The proof is
Theorem~\textup{\ref{thm:heisenberg-one-particle-sewing}}.
\end{theorem}

The Fredholm determinant formula is the Gaussian case of the sewing
programme: because the Heisenberg has only a double-pole OPE, the
multi-particle sewing operator factorizes into a product of
one-particle operators, and the trace over the Fock space reduces to
a determinant.  For non-Gaussian algebras (Kac--Moody, Virasoro,
$\mathcal{W}$-algebras), the sewing operator does not factorize, and
the partition function requires the full HS-sewing machinery of
Theorem~\ref{thm:thqg-intro-hs-general}.

\subsection{The $\hat{A}$-genus and index transgression}
\label{subsec:thqg-intro-ahat-transgression}
\index{A-hat genus@$\hat{A}$-genus!appearance in modular characteristic}
\index{transgression!of $\hat{A}$-genus}

The generating function of the genus-$g$ obstruction classes is the
$\hat{A}$-genus
(Theorem~\ref{thm:modular-characteristic}):
\begin{equation}\label{eq:thqg-intro-ahat-generating}
\sum_{g \ge 1}
\mathrm{obs}_g(\cA) \cdot \hbar^{2g}
\;=\;
\kappa(\cA) \cdot \bigl(\hat{A}(i\hbar) - 1\bigr)
\;=\;
\kappa(\cA) \cdot
\left(\frac{i\hbar/2}{\sinh(i\hbar/2)} - 1\right).
\end{equation}
The $\hat{A}$-genus appears because the obstruction classes
$\mathrm{obs}_g = \kappa \cdot \lambda_g$ (proved for uniform-weight algebras at all genera, and unconditionally at $g=1$) are the Chern characters
of the Hodge bundle, and $\operatorname{td}(\mathbb{E}^\vee)
= \hat{A}(i\hbar)$ by the Grothendieck--Riemann--Roch theorem applied
to the universal curve.

The index transgression is the passage from characteristic classes on
$\overline{\mathcal{M}}_g$ to partition functions on curves: the
$\hat{A}$-genus of the Hodge bundle transgresses to the modular
discriminant $\Delta(\tau)$ on the moduli of elliptic curves.  For
the Heisenberg, this gives
$Z_1(\cH_k) = \eta(\tau)^{-k}$, where $\eta(\tau) = q^{1/24}
\prod_{n \ge 1}(1 - q^n)$ is the Dedekind eta function.
For the Virasoro at central charge~$c$, the genus-$1$ partition
function is $Z_1(\mathrm{Vir}_c) = q^{-c/24}\,\prod_{n \ge 2}
(1-q^n)^{-1}$, and the one-loop effective action is
$\log Z_1 = -\frac{c}{24}\log q + \sum_{n \ge 2} \log(1-q^n)^{-1}$.

\subsection{The critical string at $c = 26$}
\label{subsec:thqg-intro-critical-string}
\index{critical string!$c = 26$}
\index{bosonic string!from Koszul duality}

The bosonic string critical dimension emerges from Koszul duality
alone.  The Virasoro algebra $\mathrm{Vir}_c$ is Koszul dual to
$\mathrm{Vir}_{26-c}$
(Chapter~\ref{chap:w-algebras}), and the
complementarity sum is
\[
c + c' = c + (26 - c) = 26.
\]
The anomaly-cancellation condition for the string is
$c_{\mathrm{matter}} + c_{\mathrm{ghost}} = 0$, where the ghost
contribution is $c_{\mathrm{ghost}} = -26$.  From the bar-complex
perspective, the Koszul dual of the matter Virasoro is
$\mathrm{Vir}_c^! \simeq \mathrm{Vir}_{26-c}$
(Chapter~\ref{chap:w-algebras}), and the complementarity sum
$c + c' = c + (26 - c) = 26$ recovers the critical dimension.
The ghost system ($bc$ at $c = -26$) is the physical implementation
of the Koszul dual; the cobar $\Omega(\barB(\mathrm{Vir}_c))$
recovers $\mathrm{Vir}_c$ itself by bar-cobar inversion (Theorem~B),
not the dual
(Remark~\ref{rem:bosonic-string-26}).

The perturbative finiteness of the bosonic string at $c = 26$ is
the statement that the total BRST complex has $\kappa = 0$: the
genus tower is unobstructed, and the bar complex is nilpotent at
all genera.  Equivalently, the modular characteristic of the
matter-ghost system vanishes:
\[
\kappa(\mathrm{Vir}_{26} \otimes bc_{-26})
= \kappa(\mathrm{Vir}_{26}) + \kappa(bc_{-26})
= 13 + (-13) = 0,
\]
using the additivity of $\kappa$ under tensor product: the matter
Virasoro at $c = 26$ has $\kappa = 13$, while the $bc$ ghost system
at $c = -26$ has $\kappa = -13$.  The critical dimension is not an
input but a consequence of Koszul duality.

\subsection{The sewing envelope and analytic continuation}
\label{subsec:thqg-intro-sewing-envelope}
\index{sewing envelope}
\index{analytic continuation!via sewing envelope}

The algebraic bar complex is a dense skeleton of the analytic object.
The passage from algebra to analysis is the \emph{sewing envelope}
$\cA^{\mathrm{sew}}$: the Hausdorff completion of~$\cA$ in the
all-amplitude seminorm topology.  The sewing envelope carries
the structure of a topological chiral algebra, and the HS-sewing
condition ensures that the sewing amplitudes extend from the
algebraic core to the completed object.

The analytic bar coalgebra $\barB^{\mathrm{an}}(\cA)$ is the graph-norm
closure of the algebraic bar complex inside the completed tensor
algebra of the sewing envelope.  The analytic Koszul pair
$(\cA^{\mathrm{sew}}, (\cA^!)^{\mathrm{sew}})$ is the completed
version of the algebraic Koszul pair, and the analytic shadow
$Q_g^{\mathrm{an}}(\cA)$ is the genus-$g$ factorization homology of
the analytic bar coalgebra.

The platonic chain relating the algebraic and analytic objects is
\begin{equation}\label{eq:thqg-intro-platonic-chain}
\cA_{\mathrm{alg}} \;\subset\;
\cA^{\mathrm{sew}} \;\rightsquigarrow\;
\mathcal{F}_\cA \in \mathrm{IndHilb} \;\rightsquigarrow\;
Q_\bullet^{\mathrm{an}}(\cA),
\end{equation}
where $\cA_{\mathrm{alg}}$ is the algebraic core,
$\cA^{\mathrm{sew}}$ is the sewing envelope,
$\mathcal{F}_\cA$ is the factorization algebra valued in ind-Hilbert
spaces (in the sense of Moriwaki~\cite{Moriwaki26a}), and
$Q_\bullet^{\mathrm{an}}(\cA)$ is the sequence of analytic shadows.
The Heisenberg sewing theorem
(Theorem~\ref{thm:thqg-intro-heisenberg-sewing})
is the first completed step in this chain.  The remaining steps are
the subject of the analytic sewing programme.


% ======================================================================
%  SECTION 8: GUIDE TO THE DEVELOPMENT
% ======================================================================

\section{Guide to the development}
\label{sec:thqg-intro-guide}
\index{guide to the development}

The results of this monograph are organized along ten principal lines,
each building on the categorical logarithm of Part~\ref{part:bar-complex},
the nonlinear engine of Part~\ref{part:characteristic-datum}, and the $E_1$ wing of Part~\ref{part:standard-landscape},
then extending outward through the standard landscape (Part~\ref{part:physics-bridges}),
the bridges to mathematical physics (Part~\ref{part:seven-faces}), and the frontier
(Part~VII).  The
holographic datum $\mathcal{H}(\cA)$ of
\S\ref{subsec:holographic-datum-sixfold} provides the unifying thread:
each result is a projection or extraction from the single MC
element~$\Theta_\cA$.

\subsection{The algebraic engine}
\label{subsec:thqg-intro-guide-engine}

The algebraic engine consists of Theorems A--D and H, proved in
Part~\ref{part:bar-complex}.  The constructions are:
\begin{enumerate}[label=\textup{(\roman*)}]
\item The bar-cobar adjunction $\barB_X \dashv \Omega_X$ with
  Verdier intertwining (Theorem~A,
  Chapter~\ref{chap:bar-cobar-adjunction}).
\item Bar-cobar inversion on the Koszul locus (Theorem~B,
  Chapter~\ref{chap:bar-cobar-adjunction}).
\item The Lagrangian complementarity decomposition
  $\mathbf{Q}_g(\cA) \oplus \mathbf{Q}_g(\cA^!)$
  (Theorem~C, Chapter~\ref{chap:higher-genus}).
\item The scalar modular characteristic $\kappa(\cA)$ with
  $\hat{A}$-genus generating function (Theorem~D,
  Chapter~\ref{chap:higher-genus}).
\item Polynomial chiral Hochschild cohomology on the Koszul locus
  (Theorem~H, Chapter~\ref{chap:deformation-theory}).
\end{enumerate}
These five theorems are projections of the single identity
$D_\cA^2 = 0$ at the convolution level.

\subsection{The shadow obstruction tower}
\label{subsec:thqg-intro-guide-shadow}

The nonlinear characteristic layer (Part~\ref{part:bar-complex},
Chapter~\ref{app:nonlinear-modular-shadows}) extends the scalar
invariant $\kappa$ to the full MC element $\Theta_\cA$:
\begin{enumerate}[label=\textup{(\roman*)}]
\item Construction of the modular convolution dg Lie algebra
  $\gAmod$ and its cyclic subcomplex $\Defcyc^{\mathrm{mod}}(\cA)$
  (Chapter~\ref{chap:higher-genus}).
\item The bar-intrinsic MC element
  $\Theta_\cA := D_\cA - d_\cA^{(0)}$
  (Theorem~\ref{thm:mc2-bar-intrinsic}).
\item Finite-order projections: $\kappa$ (arity~$2$),
  $\mathfrak{C}$ (arity~$3$), $\mathfrak{Q}$ (arity~$4$).
\item Shadow archetypes: G/L/C/M classification of the standard
  landscape.
\item Clutching laws for the quartic class
  (Theorem~\ref{thm:nms-clutching-law-modular-resonance}).
\item Shadow algebra $\cA^{\mathrm{sh}} = H_\bullet(\Defcyc^{\mathrm{mod}}(\cA))$
  as universal home for modular invariants.
\end{enumerate}

\subsection{The standard landscape}
\label{subsec:thqg-intro-guide-landscape}

Part~\ref{part:physics-bridges} computes the algebraic engine on every major family.  Each
chapter produces the bar complex, Koszul dual, genus tower, modular
characteristic, and shadow archetype.  The families are:
Heisenberg, free fermion, lattice, $\beta\gamma$, symplectic fermion,
affine Kac--Moody (all types), principal $\mathcal{W}$-algebras
(all types), Virasoro, $\mathcal{W}_N$ ($N \ge 3$), deformation
quantization families, Yangians, and toroidal/elliptic algebras.

\subsection{The holographic triangle}
\label{subsec:thqg-intro-guide-triangle}

The holographic programme
(\S\ref{sec:thqg-intro-bar-to-holographic}) assembles the engine
into the six-fold datum $\mathcal{H}(\cA)$.  The Swiss-cheese
factorization (Chapter~\ref{sec:swiss-cheese}), the
bulk-boundary-line triangle, the collision filtration
(\S\ref{sec:thqg-intro-collision-yang-baxter}), and the four recovery
theorems (\S\ref{subsec:thqg-intro-four-recovery}) produce a rigid
package whose internal consistency is guaranteed by the MC equation.

\subsection{The shadow connection}
\label{subsec:thqg-intro-guide-connection}

The shadow connection
(\S\ref{sec:thqg-intro-shadow-connection}) gives the general flat
modular connection on the bar-cohomological / derived-coinvariant
package.  On the affine genus-$0$ comparison surface it identifies with
the KZ connection, and on the Virasoro degenerate-module surface its
higher collision residues yield the BPZ equations.  The genus-$1$
extension involves the
Hodge bundle and the modular characteristic~$\kappa$.

\subsection{The Yangian extraction}
\label{subsec:thqg-intro-guide-yangian}

The collision-to-Yang--Baxter programme
(\S\ref{sec:thqg-intro-collision-yang-baxter}) constructs the
dg-shifted Yangian from the bar complex.  The CYBE is the Arnold
relation at codimension~$2$.  The modular Yangian carries the genus
corrections as a pro-MC element.  The Drinfeld--Kohno theorem is
the affine comparison package from KZ monodromy to the
corresponding quantum-group side
(Chapter~\ref{chap:yangians}).

\subsection{Perturbative finiteness}
\label{subsec:thqg-intro-guide-finiteness}

The analytic completion programme
(\S\ref{sec:thqg-intro-perturbative-finiteness}) upgrades the
algebraic bar complex to convergent partition functions.  The
HS-sewing criterion (Theorem~\ref{thm:general-hs-sewing}) covers
the standard landscape.  The Heisenberg sewing theorem is the
Gaussian seed.  The critical string at $c = 26$ is an output of
the Virasoro complementarity sum.

\subsection{Completion and the MC frontier}
\label{subsec:thqg-intro-guide-completion}

The MC frontier (\S\ref{sec:mc-frontier-intro}) records the boundary
between proved and open territory.  MC1 through MC4 are proved; MC5
is partially proved (analytic HS-sewing package at all genera;
genuswise BV/BRST/bar identification conjectural).
MC3 is proved for all simple types on the evaluation-generated core
(Theorem~\ref{thm:categorical-cg-all-types},
Corollary~\ref{cor:mc3-all-types}); the residual DK-4/5 problem
(extension and completion beyond evaluation modules) is downstream
of MC3.  The remaining
frontiers are non-principal W-algebra duality, analytic completion, and the
genuswise BV/BRST/bar comparison.  The completed bar-cobar
theory (MC4, Theorem~\ref{thm:completed-bar-cobar-strong}) gives a
completion-closed homotopy equivalence.  The MC4 splitting (positive
towers solved, resonant towers reduced to finite problems) is the
structural resolution of the completion obstacle.

\subsection{Bridges to physics}
\label{subsec:thqg-intro-guide-bridges}

Part~\ref{part:seven-faces} connects the engine to Feynman diagrams, BV-BRST quantization
(Chapter~\ref{ch:bv-brst}), Yang--Mills boundary theory,
and geometric Langlands.  The bar complex is the BRST
complex; the MC equation is the quantum master equation
$\hbar\Delta S + \frac{1}{2}\{S,S\} = 0$; the shadow connection is
the gauge connection of the BV formalism.

\subsection{Volume~II and beyond}
\label{subsec:thqg-intro-guide-vol2}

Volume~II applies the engine to the full bulk/boundary/line triangle
in~$3$d holomorphic-topological QFT\@.  The bar complex's differential
is $C$-direction factorization; its coproduct is $R$-direction
factorization; together they form the Swiss-cheese coalgebra on
$\mathrm{FM}_k(\mathbb{C}) \times \operatorname{Conf}_k(\mathbb{R})$.
The recognition theorem, homotopy-Koszulity of
$\mathrm{SC}^{\mathrm{ch,top}}$, PVA descent, and the derived
conformal-block line-operator identification are proved in Volume~II.
The PVA quantization bridge and holographic applications complete
the programme.

\subsection{The Koszulness characterization programme}
\label{subsec:thqg-intro-guide-koszulness}

The meta-theorem on chiral Koszulness
(Theorem~\ref{thm:koszul-equivalences-meta}) lists twelve
characterizations of the Koszul property for chiral algebras:
items~(i)--(x) are ten \emph{unconditional} equivalences;
item~(xi), the Lagrangian criterion, is \emph{conditional}
on perfectness and nondegeneracy hypotheses;
item~(xii), D-module purity, is \emph{one-directional}
((xii)~$\Rightarrow$~(x) proved, converse open).
The conditions range across several mathematical
disciplines:
\begin{enumerate}[label=\textup{(\roman*)}]
\item $A_\infty$-formality of the bar cohomology.
\item Shadow-formality: the shadow obstruction tower is equivalent to the
  $L_\infty$-formality obstruction tower at arities $2$, $3$, $4$.
\item $E_2$-formality of chiral Hochschild cohomology.
\item Genus-$0$ curve independence: Koszulness on~$\mathbb{P}^1$ implies
  Koszulness on any smooth curve.
\item Graph-sum truncation: the graph sum for $\Theta_\cA^{\le r}$
  terminates at finitely many graphs.
\item Universal Koszulness: $V_k(\mathfrak{g})$ is chirally Koszul for
  all~$k$.
\item PBW universality: freely strongly generated vertex algebras are
  chirally Koszul.
\item Barr--Beck--Lurie: the monadic characterization via the
  bar-cobar monad.
\item Factorization homology concentration: $\int_X \cA$ is concentrated
  in the expected cohomological degrees.
\item FM boundary acyclicity: the restriction of the bar complex to
  the boundary of $\overline{C}_n(X)$ is acyclic.
\item Tropical Koszulness: the tropicalization of the bar complex is
  acyclic.
\item Lagrangian criterion: the Lagrangian complement
  $\mathbf{Q}_g(\cA^!)$ is concentrated in the expected degrees.
  \emph{(Conditional on perfectness and nondegeneracy hypotheses.)}
\end{enumerate}
The ten unconditional equivalences are proved by chain-level
arguments connecting each condition to the degeneration of the bar
spectral sequence at~$E_2$.  The Lagrangian criterion requires
additional hypotheses on the ambient pairing (see
Theorem~\ref{conj:lagrangian-koszulness}).
The programme is recorded in
\S\ref{sec:concordance-koszulness-programme} of the concordance.

\subsection{Summary of dependencies}
\label{subsec:thqg-intro-guide-dependencies}

The logical dependencies among the principal constructions are:
\begin{center}
\small
\begin{tabular}{ll}
\textbf{Construction} & \textbf{Depends on} \\ \hline
Bar complex $\barB_X(\cA)$ & chiral algebra $\cA$, FM compactification \\[2pt]
Koszul dual $\cA^!$ & Theorem~A (Verdier intertwining) \\[2pt]
MC element $\Theta_\cA$ & $D_\cA^2 = 0$ (formal consequence of $\partial^2 = 0$ on $\overline{\mathcal{M}}_{g,n}$) \\[2pt]
Shadow obstruction tower & projections of $\Theta_\cA$ \\[2pt]
$r$-matrix $r(z)$ & collision filtration on $\barB_X(\cA)$ \\[2pt]
CYBE & Arnold relation on $\overline{C}_3(X)$ \\[2pt]
Yangian $\mathcal{Y}^{\mathrm{dg}}_T$ & $r$-matrix + RTT equation \\[2pt]
Shadow connection $\nabla^{\mathrm{hol}}$ & $\operatorname{Sh}_{g,n}(\Theta_\cA)$ \\[2pt]
 KZ connection (affine comparison surface) &
   $\nabla^{\mathrm{hol}}|_{g=0}$ \\[2pt]
 BPZ equations (Virasoro degenerate-module surface) &
   null vectors + shadow connection \\[2pt]
HS-sewing & OPE growth bounds \\[2pt]
$c = 26$ & Virasoro Koszul duality + $\kappa_{\mathrm{tot}} = 0$
\end{tabular}
\end{center}
Every entry in the right column is either a geometric input (FM
compactification, boundary stratification of moduli) or a consequence
of the identity $D_\cA^2 = 0$.  The holographic programme is a tree
rooted at this single identity.

\subsection{The master invariant table}
\label{subsec:thqg-intro-guide-master-table}

The following table collects the principal numerical invariants
produced by the holographic datum for the standard families.  Each
entry is a projection of~$\Theta_\cA$.

\begin{center}
\small
\renewcommand{\arraystretch}{1.15}
\begin{tabular}{lcccc}
\textbf{Family}
  & $\kappa$ & $c + c'$ & $r_{\max}$ & \textbf{Class} \\ \hline
Heisenberg $\cH_k$ & $k$ & $2$ & $2$ & G \\[2pt]
Free fermion $\psi_k$ & $k/4$ & $1$ & $2$ & G \\[2pt]
Lattice $V_\Lambda$ & $\operatorname{rank}(\Lambda)$ & $2\operatorname{rank}(\Lambda)$ & $2$ & G \\[2pt]
$\beta\gamma$ & $-1$ & $-2$ & $4$ & C \\[2pt]
$\widehat{\mathfrak{sl}}_2{}_k$
  & $\frac{3(k+2)}{4}$
  & $6$ & $3$ & L \\[2pt]
$\widehat{\mathfrak{g}}_k$
  & $\frac{(k+h^\vee)\dim\mathfrak{g}}{2h^\vee}$
  & $2\dim\mathfrak{g}$ & $3$ & L \\[2pt]
$\mathrm{Vir}_c$ & $c/2$ & $26$ & $\infty$ & M \\[2pt]
$\mathcal{W}^k(\mathfrak{sl}_N)$
  & $c\,(H_N - 1)$
  & $2(N{-}1)(2N^2{+}2N{+}1)$
  & $\infty$ & M \\[2pt]
$Y(\mathfrak{sl}_N)$ (Yangian)
  & $N^2 - 1$ & $2(N^2-1)$ & $3$ & L
\end{tabular}
\end{center}

\noindent
The modular characteristic $\kappa$ determines the one-loop physics;
the conductor $c + c'$ determines the anomaly cancellation; the
shadow depth $r_{\max}$ determines the gravitational complexity.  All
three invariants are computable from the bar complex of~$\cA$, and
all three are projections of the holographic datum~$\mathcal{H}(\cA)$.

The seven sections of this supplement, from the Swiss-cheese
factorization to the master invariant table, form a self-contained
executive summary of the holographic content of the monograph.  The
full proofs and computations appear in Parts~\ref{part:bar-complex}--\ref{part:physics-bridges} of Volume~I and in
Volume~II.
